\documentclass[10pt]{ctexart}
\usepackage{graphicx}     % 图片插入
\usepackage{xcolor}
\usepackage{textcomp}
\usepackage{url}
\usepackage{listings}
% a4版本
%\usepackage[margin=1.2in]{geometry}
% 下面是a5版本
\usepackage{geometry}     % 设置页边距
% 设置 A5 纸张尺寸和 1 英寸的边距
\geometry{a5paper, left=0.5in, right=0.5in, top=0.5in, bottom=0.5in}


\definecolor{darkgray}{rgb}{.4,.4,.4}

\lstdefinelanguage{JavaScript}{
  keywords={typeof, new, true, false, catch, function, return, null, catch, switch, var, if, in, while, do, else, case, break},
  ndkeywords={class, export, boolean, throw, implements, import, this},
  sensitive=false,
  comment=[l]{//},
  morecomment=[s]{/*}{*/},
  morestring=[b]',
  morestring=[b]"
}

\lstset{
    aboveskip={.3\baselineskip},
    basicstyle=\scriptsize\ttfamily\linespread{4},
    breaklines=false,
    columns=flexible,
    commentstyle=\color[rgb]{0.127,0.427,0.514}\ttfamily\itshape,
    escapechar=@,
    extendedchars=true,
    frame=single,
    identifierstyle=\color{black},
    inputencoding=latin1,
    keywordstyle=\color[HTML]{228B22}\bfseries,
    language=JavaScript,
    ndkeywordstyle=\color[HTML]{228B22}\bfseries,
    numbers=left,
    numberstyle=\tiny,
    prebreak = \raisebox{0ex}[0ex][0ex]{\ensuremath{\hookleftarrow}},
    stringstyle=\color[rgb]{0.639,0.082,0.082}\ttfamily,
    upquote=true,
    showstringspaces=false,
}

\lstset{literate=%
   *{0}{{{\color{darkgray}0}}}1
    {1}{{{\color{darkgray}1}}}1
    {2}{{{\color{darkgray}2}}}1
    {3}{{{\color{darkgray}3}}}1
    {4}{{{\color{darkgray}4}}}1
    {5}{{{\color{darkgray}5}}}1
    {6}{{{\color{darkgray}6}}}1
    {7}{{{\color{darkgray}7}}}1
    {8}{{{\color{darkgray}8}}}1
    {9}{{{\color{darkgray}9}}}1
  }


 

  \begin{document}

\newpage
% 生成目录
\tableofcontents
\newpage

\clearpage % 强制所有浮动体（图片、表格等）输出并分页 避免图片跨页：使用 \clearpage 或 \newpage 来强制图片排版在当前页面。
% \newpage 
\section{Docker核心概念总结}
{\footnotesize % 脚注大小
Docker概述：

Docker是一个开源的容器化平台，用于自动化应用程序的部署、扩展和管理。它通过将应用程序及其所有依赖项打包成一个容器，确保它们在任何环境中都能一致地运行。

容器化：Docker容器是轻量级的虚拟化技术，允许在一个宿主操作系统上运行多个隔离的应用。

Docker的组成：

Docker Engine：包括客户端和服务器，负责容器的构建和管理。

Docker镜像：包含应用运行所需的所有文件、库、依赖等，是容器的可执行包。

Docker容器：镜像的运行实例，是可以启动、停止、删除和管理的应用环境。

Docker Hub：Docker的公共镜像仓库，用户可以从中获取共享的镜像。

常用Docker命令：

docker run <image>：从指定镜像创建并启动一个新的容器。

docker ps：查看当前运行中的容器。

docker stop <container>：停止正在运行的容器。

docker rm <container>：删除停止的容器。

docker pull <image>：从Docker Hub或指定仓库拉取镜像。

docker build -t <image-name> <dockerfile>：使用Dockerfile构建镜像。

docker exec -it <container> /bin/bash：进入容器进行交互式操作。

docker logs <container>：查看容器的日志输出。

docker images：查看本地存储的所有镜像。

Docker容器生命周期：

创建：使用docker run或docker create命令从镜像创建容器。

启动：容器启动后，可以使用docker start命令来启动一个已创建的容器。

停止：使用docker stop停止正在运行的容器。

删除：使用docker rm删除已停止的容器。

Docker网络：

Bridge网络：默认网络模式，用于容器间的通信，容器通过虚拟网桥互联。

Host网络：容器与主机共享网络栈，直接使用主机的IP和端口。

None网络：不为容器分配网络，容器无法与外界通信。

Docker的存储：

Volumes：Docker容器中的数据存储。数据保存在宿主机上，且容器重启时数据不会丢失。

Bind Mounts：将宿主机的目录挂载到容器内，容器的修改会反映到宿主机上。

Tmpfs：将数据存储在内存中，当容器停止时数据会丢失。

Docker的安全性：用户权限：Docker容器以root用户身份运行，但可以通过USER指令指定非root用户运行。安全扫描：使用docker scan命令扫描镜像中的安全漏洞。隔离：通过容器的命名空间和cgroups技术来实现容器之间的隔离。
}


\clearpage
% \newpage
\section{常用Kubernetes命令总结}
{\footnotesize % 脚注大小
集群管理命令：

kubectl get nodes：查看集群中的所有节点。

kubectl get pods：查看所有Pod的状态。

kubectl get services：查看所有Service。

kubectl get deployments：查看所有Deployment。

Pod管理命令：

kubectl create -f pod.yaml：通过YAML文件创建Pod。

kubectl describe pod <pod-name>：查看指定Pod的详细信息。

kubectl delete pod <pod-name>：删除指定Pod。

kubectl logs <pod-name>：查看Pod中的容器日志。

Deployment管理命令：

kubectl create -f deployment.yaml：创建Deployment。

kubectl get deployments：查看所有Deployment。

kubectl set image deployment/<deployment-name> <container-name>=<image>：更新Deployment的镜像。

kubectl scale deployment <deployment-name> --replicas=<number>：调整Deployment副本数。

kubectl rollout status deployment/<deployment-name>：查看Deployment的滚动更新状态。

kubectl rollout undo deployment/<deployment-name>：回滚Deployment。

Service管理命令：

kubectl expose pod <pod-name> --port=<port>：为Pod暴露Service。

kubectl expose deployment <deployment-name> --type=LoadBalancer --port=<port>：为Deployment暴露LoadBalancer类型的Service。

kubectl get svc：查看所有Service的状态。

调度命令：

kubectl get pod <pod-name> -o wide：查看Pod的详细信息，包括所在节点。

kubectl logs <pod-name> -c <container-name>：查看指定容器的日志。

集群调试命令：

kubectl describe pod <pod-name>：获取Pod的详细描述信息。

kubectl get events：查看集群的事件，帮助排查问题。

kubectl exec -it <pod-name> -- /bin/bash：进入Pod中的容器进行交互。
}


\newpage
 
\title{用  \texttt{listings} 制作类 minted 风格}

\maketitle

不少用户不会配置 minted 环境，我们也制作了视频教程教大家如何配置，可以到这里看看视频配置。
\url{https://www.bilibili.com/video/BV1sT4y1A7KR}


如果还是不会配置，今天给大家一个用
listings 宏包定制的类 minted 风格，这样就不用配置环境，也可以体验到 minted 样式了。



显示效果如下：
\begin{lstlisting}
var loginLayer=(function(){
    var div=document.createElement("div");
    div.innerHTML="windows1";
    div.style.display="none";
    document.body.appendChild(div);
    return div;
})();
/*document.getElementById('loginBtn').onclick=function(){
    loginLayer.style.display="block";
};*/
alert(loginLayer.style.display);

    public class HelloWorld {
        public static void main(String[] args) {
            System.out.println("Hello, World!");
        }
    }
\end{lstlisting}

\newpage
N.53 最大子数组和的问题
\begin{lstlisting}
public int maxSubArray(int[] nums) {
        int max = nums[0];
        for(int i=0;i<nums.length;i++){
            int tmpMax = nums[i];
            if(tmpMax>max) {
                max = tmpMax;
            }
         for(int j=i+1;j<nums.length;j++) {
            tmpMax += nums[j];
            if(tmpMax>max) {
                max = tmpMax;
            }
         }
        }
        System.out.println("max--"+max);
        return max;
    }
\end{lstlisting}



\newpage
\begin{figure}
    \centering
    \includegraphics[width=0.5\textwidth]{9.png} % 确保替换为你的图片路径
    % \includegraphics[width=\textwidth, height=6cm]{example.jpg}  % 设置最大宽度为页面宽度，并限制高度
    %\includegraphics[width=\textwidth, height=12cm]{9.png}  % 设置最大宽度为页面宽度，并限制高度
    \caption{示例图片}
    \label{fig:example}
\end{figure}


\clearpage  % 所有浮动体（如图片）被处理并开始新的一页


% 插入第一张图片，上方，靠左对齐，去掉标题和空白
\begin{figure}[t!]
    \vspace{-50pt}  % 向上移动图片
    \hspace{-20pt}  % 向左移动图片
    \raggedright  % 使图片靠左
    \includegraphics[height=9cm]{9.png}  % 调整图片高度
    %\caption*{}  % 无标题显示
    \vspace{-70pt}  % 减少图片下方的空白
\end{figure}

% 插入第二张图片，下方，靠左对齐，去掉标题和空白
\begin{figure}[h!]
    \raggedright  % 使图片靠左
    \includegraphics[height=9cm]{9.png}  % 调整图片高度
    %\caption*{}  % 无标题显示
    \vspace{-10pt}  % 减少图片下方的空白
\end{figure}

\clearpage

% \newpage
核心定制代码如下：
\small
\begin{verbatim}
\lstset{
    aboveskip={.3\baselineskip},
    basicstyle=\scriptsize\ttfamily\linespread{4},
    breaklines=false,
    columns=flexible,
    commentstyle=\color[rgb]{0.127,0.427,0.514}\ttfamily\itshape,
    escapechar=@,
    extendedchars=true,
    frame=single,
    identifierstyle=\color{black},
    inputencoding=latin1,
    keywordstyle=\color[HTML]{228B22}\bfseries,
    language=JavaScript,
    ndkeywordstyle=\color[HTML]{228B22}\bfseries,
    numbers=left,
    numberstyle=\tiny,
    prebreak = \raisebox{0ex}[0ex][0ex]{\ensuremath{\hookleftarrow}},
    stringstyle=\color[rgb]{0.639,0.082,0.082}\ttfamily,
    upquote=true,
    showstringspaces=false,
}
\end{verbatim}

% \lstinputlisting{filename.js}







\end{document}
